\SourceLeaf{245}%
tione retinebimus, in qua in art. praec. acceptae fuerunt, modo indefinite ad punctum indefinitum cuiuslibet linearum brevissimarum referantur.

Lineae brevissimae omnes, quae sunt aequalis longitudinis \(r\), terminabuntur ad aliam lineam, cuius longitudinem ab initio arbitrario numeratam denotamus per \(v\). Considerari poterit itaque \(v\) tamquam functio indeterminatarum \(r,\varphi\), et si per \(\lambda'\) designamus punctum in superficie sphaerica respondens directioni elementi \(\mathrm{d}v\), nec non per \(\xi',\eta',\zeta'\) coordinatas huius puncti respectu centri sphaerae, habebimus:
\[
\frac{\mathrm{d}x}{\mathrm{d}\varphi}=\xi'\mathord{.}\frac{\mathrm{d}v}{\mathrm{d}\varphi},\quad
\frac{\mathrm{d}y}{\mathrm{d}\varphi}=\eta'\mathord{.}\frac{\mathrm{d}v}{\mathrm{d}\varphi},\quad
\frac{\mathrm{d}z}{\mathrm{d}\varphi}=\zeta'\mathord{.}\frac{\mathrm{d}v}{\mathrm{d}\varphi}
\]
Hinc et ex
\[
\frac{\mathrm{d}x}{\mathrm{d}r}=\xi,\quad\frac{\mathrm{d}y}{\mathrm{d}r}=\eta,\quad\frac{\mathrm{d}z}{\mathrm{d}r}=\zeta
\]
sequitur
\[
\begin{aligned}
\frac{\mathrm{d}x}{\mathrm{d}r}\mathord{.}\frac{\mathrm{d}x}{\mathrm{d}\varphi}+\frac{\mathrm{d}y}{\mathrm{d}r}\mathord{.}\frac{\mathrm{d}y}{\mathrm{d}\varphi}+\frac{\mathrm{d}z}{\mathrm{d}r}\mathord{.}\frac{\mathrm{d}z}{\mathrm{d}\varphi}
&=(\xi\xi'+\eta\eta'+\zeta\zeta')\mathord{.}\frac{\mathrm{d}v}{\mathrm{d}\varphi}\\
&=\cos\lambda\lambda'\mathord{.}\frac{\mathrm{d}v}{\mathrm{d}\varphi}
\end{aligned}
\]
Membrum primum huius aequationis, quod etiam erit functio ipsarum \(r,\varphi\), per \(S\) denotamus; cuius differentiatio secundum \(r\) suppeditat:
\[
\begin{aligned}
\frac{\mathrm{d}S}{\mathrm{d}r}&=\frac{\mathrm{d}\mathrm{d}x}{\mathrm{d}r^2}\mathord{.}\frac{\mathrm{d}x}{\mathrm{d}\varphi}+\frac{\mathrm{d}\mathrm{d}y}{\mathrm{d}r^2}\mathord{.}\frac{\mathrm{d}y}{\mathrm{d}\varphi}+\frac{\mathrm{d}\mathrm{d}z}{\mathrm{d}r^2}\mathord{.}\frac{\mathrm{d}z}{\mathrm{d}\varphi}\\
&\quad+\tfrac{1}{2}\mathord{.}\frac{\mathrm{d}((\tfrac{\mathrm{d}x}{\mathrm{d}r})^2+(\tfrac{\mathrm{d}y}{\mathrm{d}r})^2+(\tfrac{\mathrm{d}z}{\mathrm{d}r})^2)}{\mathrm{d}\varphi}\\
&=\frac{\mathrm{d}\xi}{\mathrm{d}r}\mathord{.}\frac{\mathrm{d}x}{\mathrm{d}\varphi}+\frac{\mathrm{d}\eta}{\mathrm{d}r}\mathord{.}\frac{\mathrm{d}y}{\mathrm{d}\varphi}+\frac{\mathrm{d}\zeta}{\mathrm{d}r}\mathord{.}\frac{\mathrm{d}z}{\mathrm{d}\varphi}\\
&\quad+\tfrac{1}{2}\mathord{.}\frac{\mathrm{d}(\xi\xi+\eta\eta+\zeta\zeta)}{\mathrm{d}\varphi}
\end{aligned}
\]
Sed \(\xi\xi+\eta\eta+\zeta\zeta=1\), adeoque ipsius differentiale \(=0\); et per art. praec. habemus, si etiam hic \(\rho\) denotat radium curvaturae in linea \(r\),
\[
\frac{\mathrm{d}\xi}{\mathrm{d}r}=\frac{X}{\rho},\quad\frac{\mathrm{d}\eta}{\mathrm{d}r}=\frac{Y}{\rho},\quad\frac{\mathrm{d}\zeta}{\mathrm{d}r}=\frac{Z}{\rho}
\]
Ita obtinemus
\[
\frac{\mathrm{d}S}{\mathrm{d}r}=\frac{1}{\rho}\mathord{.}(X\xi'+Y\eta'+Z\zeta')\mathord{.}\frac{\mathrm{d}v}{\mathrm{d}\varphi}=\frac{1}{\rho}\mathord{.}\cos L\lambda'\mathord{.}\frac{\mathrm{d}v}{\mathrm{d}\varphi}=0
\]
quoniam manifesto \(\lambda'\) iacet in circulo maximo, cuius polus \(L\). Hinc itaque concludimus, \(S\) independentem esse ab \(r\) et proin functionem solius \(\varphi\). At pro \(r=0\) manifesto fit \(v=0\), et proin etiam \(\frac{\mathrm{d}v}{\mathrm{d}\varphi}=0\), nec non \(S=0\), independenter a \(\varphi\). Necessario itaque generaliter esse debebit \(S=0\), adeoque \(\cos\lambda\lambda'=0\), i. e. \(\lambda\lambda'=90^0\). Hinc colligimus
